\documentclass[10pt,a4paper]{amsart}
\usepackage[margin=19mm]{geometry}
\usepackage{amsmath,amssymb,mathtools}
\usepackage[scr=rsfs,cal=euler]{mathalfa}
\usepackage{xfrac}
\usepackage[unicode,hidelinks,bookmarks=false]{hyperref}
\newcommand{\ie}{i.e.\ }
\newcommand{\cf}{cf.\ }
\newcommand{\resp}{respectively\ }
\newcommand{\Subsection}{\subsection}
\providecommand{\nobreakdash}{\nobreak\hbox{-}}
\setlength{\emergencystretch}{2em}
\multlinegap0pt
\usepackage{bm}
\usepackage{bbm}
\usepackage{dsfont}
\usepackage{xspace}
\usepackage{cancel}
\usepackage{mathtools}
\usepackage{amsthm}
\makeatletter
\@ifundefined{theorem}{\newtheorem{theorem}{Theorem}}{}
\@ifundefined{lemma}{\newtheorem{lemma}[theorem]{Lemma}}{}
\makeatother
\DeclareMathOperator{\Gal}{Gal}
\DeclareMathOperator{\rk}{rk}
\newcommand{\CC}{\mathbb{C}}
\newcommand{\RR}{\mathbb{R}}
\newcommand{\fo}{\mathfrak{o}}
\title[Integral Diophantine approximation on varieties]{Integral Diophantine approximation on varieties}
\author{Zhizhong Huang, Florian Wilsch}
\date{Source: arXiv:2509.03998v2 (CC BY 4.0)}
\begin{document}
\maketitle

\noindent\emph{Source and licence.} Integral Diophantine approximation on varieties. Version arXiv:2509.03998v2 (CC BY 4.0), distributed under the Creative Commons Attribution 4.0 International licence, as recorded in the arXiv licence metadata for this exact version and shown on the abstract page https://arxiv.org/abs/2509.03998v2. Full rights notice: licenses/arxiv-2509.03998-CC-BY-4.0.txt.\par\smallskip

\noindent\textit{Editorial scope.} Counted unit: the Lemma whose source label is \texttt{lem:units-of-norm-1} in the article (source0.tex lines 670--672, source label \texttt{lem:units-of-norm-1}), delivered together with the article's own complete original proof of it (source0.tex lines 673--707, one \texttt{proof} environment of 35 lines). No mathematics is written, repaired or re-derived: statement and proof are the article's own text line for line. Required context retained uncounted: eq:norm-1-rank (lines 442--444); lem:geom-numbers (lines 663--665). Every definition, display and label of those lines is retained unchanged. Omitted from this package and not redistributed: 1--441, 445--662, 666--669, 708--1482. Nothing inside a retained excerpt is omitted or altered: every retained body is delivered exactly as the article's lines are written, so no omission receipt of any kind accompanies this package. Citations: the retained excerpts carry no citation command at all, so no bibliography entry of the article is transcribed and no reference is redistributed. Reference adaptation: none was needed and none was made; every reference inside the delivered unit resolves inside it, so no unresolved reference or citation is produced. Printed slips kept literal: no printed word, symbol, hypothesis or number of the counted statement or of its proof was altered. Layout and class: the article's own class is \texttt{amsart} with packages including fontenc, inputenc, babel, lmodern, amssymb, amsmath, amsthm, amsfonts, mathtools, dsfont, hyperref, natbib, mathrsfs, enumitem; the wrapper is the lane's pinned \texttt{amsart} editorial wrapper at 10pt on A4, which restates the theorem-like environments the retained text needs and adds the article's own macro definitions that the retained text needs (\texttt{\textbackslash CC}, \texttt{\textbackslash Gal}, \texttt{\textbackslash RR}, \texttt{\textbackslash fo}, \texttt{\textbackslash rk}), with extra wrapper packages (bm, bbm, dsfont, xspace, cancel, mathtools). The delivered numbering is the unit's own. Assets: no retained span contains a figure environment, a graphics-inclusion command, a PDF-page inclusion, a table or a TikZ picture; no image file is copied into this package, the package declares no asset dependency and no image file is redistributed with it. Licence: version arXiv:2509.03998v2 of this article is distributed under the Creative Commons Attribution 4.0 International licence, as recorded in the arXiv licence metadata for this exact version and shown on the abstract page https://arxiv.org/abs/2509.03998v2. A full rights notice is delivered at licenses/arxiv-2509.03998-CC-BY-4.0.txt. Characters: every retained excerpt is pure ASCII, so the delivered engine, XeLaTeX with its default Latin Modern text font, reports no missing character.
	\begin{equation}\label{eq:norm-1-rank}
		\rk\fo_K^1=\#\infty_K - \#\infty_k,
	\end{equation}

\begin{lemma}\label{lem:geom-numbers}
	Let $n\ge 1$ and $\Lambda$ be a lattice in $\RR^n$ of full rank. Then for every $\delta_1,\delta_2>0$, there is a lattice point $x\in \Lambda \cap [\delta_1^{-1},\infty]\times [-\delta_2,\delta_2]^{n-1}$.
\end{lemma}

\begin{lemma}\label{lem:units-of-norm-1}
	Let $K/k$ be a quadratic extension and $U$ be a subgroup of full rank of the group $\fo_K^1$ of units of norm $1$. Let $v_0$ be a place of $k$ such that $K\subset k_v$; let $w_0$ be a place above $v_0$. For every $\delta>0$, there is a unit $\epsilon\in U$ such that $|\epsilon|_{w_0}<\delta$ and $1/(1+ \delta) < |\epsilon|_{w} < 1+\delta$ for archimedean places $w$ with $w\nmid v_0$.
\end{lemma}

\begin{proof}
	Let 
	\begin{align*}
		\phi \colon \fo_K^\times &\to \RR^{\infty_K},\\\epsilon & \mapsto (\log |\epsilon|_w)_{w\in \infty_K};
	\end{align*}
	we shall begin by describing the image of $\fo_K^1$ under this map by means of a linear subspace.
	The set of archimedean places $\infty_k$ is the union of splitting places $\infty_k^{\mathrm s}$ and inert places $\infty_k^{\mathrm i}$. 

	Let $v\in \infty_k^{\mathrm s}$ be one of the former for now, that is, a place that has two different places $w_1$ and $w_2$ lying above it; note that this case includes $v=v_0$ by assumption, and in this case, label $w_2 = w_0$.
	The two places $w_1$ and $w_2$ are swapped by the nontrivial element $\sigma$ of the Galois group $\Gal(K/k)$, hence
	\begin{equation*}
		1 = |N(\epsilon)|_{w_1} = |\epsilon|_{w_1}|\sigma(\epsilon)|_{w_1} = |\epsilon|_{w_1}|\epsilon|_{w_2}.
	\end{equation*}
	Logarithmically, this translates to
	$f_{v}(\phi(\epsilon))=0$, where
	\begin{equation*}
		f_{v} = x_{w_1}+x_{w_2}\colon \RR^{\infty_K}\to \RR
	\end{equation*}
	is the sum of the two projections.
	Let $b_{v}\in\RR^{\infty_K}$ be the vector with $1$ in the $w_1$-entry and $-1$ in the $w_2$-entry.

	Let now $v\in \infty_k^{\mathrm i}$ be an inert archimedean place and $w$ be the only place above $v$. Then $\RR\cong k_{v} \subset K_{w}\cong \CC$, and $\sigma$ acts by complex conjugation, so that
	\begin{equation*}
		1 = |N(\epsilon)|_{w} = |\epsilon\overline{\epsilon}|_w.
	\end{equation*}
	This implies $f_{v}(\phi(\epsilon))=0$, where $f_{v} = x_w\colon \RR^{\infty_K}\to \RR$ is the projection.
	
    Let $H\subset \RR^{\infty_K}$ be the linear subspace cut out by the $\# \infty_k$ linearly independent linear forms $(f_{v})_{v\in\infty_k}$. It is thus of dimension $\#\infty_K - \#\infty_k$ with basis $(b_{v})_{v\in \infty_k^{\mathrm s}}$.
    Note that $H$ contains the sublattice $\phi(\fo_K^1)$ of $\phi(\fo_K^\times)$, which is of full rank by~\eqref{eq:norm-1-rank}.
	In the basis $(b_{v})_{v\in\infty_k^{\mathrm s}}$, ordered so that $v_0\in \infty_k^{\mathrm s}$ comes first, the desired properties for the absolute values translate to
	\begin{equation*}
		\phi(\epsilon) \in [- \log \delta ,\infty] \times [-\log (1+\delta), \log(1+\delta)]^{\#\infty_k^s -1}.
	\end{equation*}
	Now the claim follows from Lemma~\ref{lem:geom-numbers}.
\end{proof}

\end{document}
